%% kapa: exam=1 round=35 year=2024 period=1 form=A subject=real_estate subject_ko=부동산학원론 seq=22 number=102
%% source: Q-Net 감정평가사 기출문제 (cst00309) artlSeq=5213229 file=raw/1008/5213229/1교시 문제지 원본.pdf
%% conversion: kordoc 4.24.51 -> tools/gen1.py (auto); review=proofread-sonnet; flags=none
\begin{question}{102}
부동산의 수요와 공급에 관한 설명으로 옳지 않은 것은? (단, 우하향하는 수요곡선과 우상향하는 공급곡선을 가정하며, 다른 조건은 동일함)
\begin{choices}
  \item 단기적으로 가격이 상승해도 부동산의 공급량이 크게 증가할 수 없기 때문에 공급이 비탄력적이다.
  \item 부동산의 공급량은 주어진 가격 수준에서 일정기간에 판매하고자 하는 최대수량이다.
  \item 용도전환 및 개발이 가능한 장기에는 공급의 탄력성이 커진다.
  \item 부동산의 수요량은 구매능력을 갖춘 수요자들이 구매하려는 수량이므로 유효수요를 의미한다.
  \item 공급의 가격탄력성이 작을수록 수요변화시 균형가격의 변동폭은 작지만 균형거래량의 변동폭은 크다.
\end{choices}
\end{question}
\answer{5}
